% !TeX root = ../main.tex

\chapter{Other Scientific Contributions}
\label{apx:contributions}

\section{Reproducing the reported numbers}
\label{apx:contributions:repro}

This section identifies the code that implements each part of the protocol and the files
that contain its outputs. All paths below are relative to this research's working
repository, published at \url{https://github.com/VitorHugoOli/PoiMtlNet}. The protocol
applies only to the final study in Chapter~\ref{ch:mobiwac}, as specified in
Chapter~\ref{ch:fundamentals}.

The public repository contains the model, the representation, the baselines, the code
that constructs the folds, and the statistical scripts and output files listed below.

\begin{description}
  \item[The fold partition.] Folds are built by
        \path{src/data/folds.py}, which calls scikit-learn's
        \texttt{StratifiedGroupKFold} with five splits, shuffling enabled, and the run's
        seed as the partition seed, grouping by user identifier and stratifying on the task
        label. Grouping by user is what makes the split user-disjoint: a user's
        check-ins fall on one side of every fold.
        The same routine builds the paired multitask folds, where both tasks share one
        user partition so that a user
        cannot reach one task's training data and the other task's validation data in
        the same fold.
  \item[The seeds.] Four seeds, 0, 1, 7, and 100, each repeat the
        five-fold experiment in full. Each seed determines a different user
        partition, so the design produces $4\times5=20$ fitted models per
        configuration. The design is recorded in
        \path{docs/reproducibility/mobiwac_v17/STATISTICAL_PROTOCOL.md}. A seed
        sets both the random initialization and the partition, so the four seeds draw four
        user partitions rather than repeating one, and the reported intervals cover both
        sources of variation. Four draws sample the variation across user splits without
        characterizing it. Within one seed both models are built from the same partition, which is what
        the paired comparison requires. The paired $t$ and TOST analyses compare the four
        per-seed means, while the registered Wilcoxon sensitivity analysis compares the 20
        matched seed-fold differences.
  \item[The reported scores.] Read at one saved model per fold, the epoch selected on
        validation, which is the joint-best convention: both tasks are read at the same
        checkpoint rather than each at its own best epoch. The post-hoc scorer is
        \path{src/tracking/score_joint_best.py}, and the per-fold values it produced for
        every joint cell reported in Chapter~\ref{ch:mobiwac} are stored in
        \path{docs/results/closing_data/v18/joint_best_perfold.json}.
  \item[The significance tests.] The superiority and non-inferiority tests are
        implemented in \path{research/reproducibility/mobiwac_v17/superiority_wilcoxon.py}
        and \path{research/reproducibility/mobiwac_v17/region_match_tost.py}. The values
        they consume, and from which every cell of Chapter~\ref{ch:mobiwac}'s results table
        is computed, are the per-fold records in
        \path{docs/results/closing_data/v18/joint_best_perfold.json} for the joint model
        and \path{docs/studies/closing_data/v18/data/v18_results.json} for the dedicated
        models. Both store one value per fold per seed, so the seed-level means, the paired
        differences, and the intervals the chapter reports are recomputable from them
        without rerunning any model.
  \item[The label-history benchmark.] Computed by
        \path{research/reproducibility/mobiwac_v17/autocorrelation_ceiling.py}; its
        machine-readable outputs are stored in
        \path{docs/results/embedding_eval/rescreen_cat/}.
\end{description}

Chapter~\ref{ch:fundamentals} introduces the metric conventions for macro-F1 and
accuracy at ten, and Chapter~\ref{ch:mobiwac} defines them again. This section does not
repeat those definitions.

Chapters~\ref{ch:cbic} and~\ref{ch:courb} used the sample-stratified protocol described
earlier in this appendix, and the implementation of Chapter~\ref{ch:courb} is stored in
a separate repository. Reproducing either chapter therefore requires following the
protocol reported in that chapter rather than the one described in this section.
